\section{Materiales y método}
\label{sec:metodo}

\subsection{Datos}
\label{sec:datos}

Usamos dos bancos públicos por parcela y no ajustamos el protocolo entre uno y otro.

\begin{figure}[tbp]
\centering
\includegraphics[width=\textwidth]{parcelas.pdf}
\caption{Tres parches del fold retenido del banco primario. Arriba, composición en color natural
de Sentinel-2 en la toma más cercana al primero de julio. Abajo, la anotación agronómica por
parcela. La unidad del producto es la parcela, y el catálogo que estudia este artículo es la lista
de nombres que la fila de abajo tiene permitido usar.}
\label{fig:parches}
\end{figure}

\begin{table}[tbp]
\caption{Los dos bancos. Los dos se describen con el mismo vector de \num{185} características
temporales, así que la réplica varía dato, año, región y familia de modelo, pero no
representación.}
\label{tab:bancos}
\centering
\small
\setlength{\tabcolsep}{4pt}
\begin{tabular}{lrr}
\toprule
 & primario (PASTIS) & réplica (BreizhCrops) \\
\midrule
Campaña & 2018--2019 & 2017 \\
Región & Francia metropolitana & Bretaña \\
Parcelas puntuadas & 16\,640 & 60\,000 \\
Clases & 18 & 9 \\
Clase mayoritaria & 6\,128 & 18\,113 \\
Clase minoritaria & 103 & 2 \\
Razón de desbalance & 60:1 & 9\,057:1 \\
Bloques & 5 espaciales (H3) & 2 regiones \\
Predictor & miembros preentrenados & gradiente potenciado \\
Partición del predictor & fold retenido & una región fuera \\
\bottomrule
\end{tabular}
\end{table}

El \textbf{banco primario} es PASTIS \cite{garnot2021pastis}. Lo obtenemos por la distribución
PASTIS-R, que añade series de radar alineadas \cite{saintefaregarnot2022multimodal}, pero
\emph{usamos solo su modalidad óptica}: todos los predictores leen Sentinel-2, así que llamar al
banco PASTIS-R nombraría un sensor que nunca tocamos. De él tomamos el fold retenido y
agregamos la anotación densa a nivel de parcela, con \num{16640} parcelas sobre \num{496} parches
de imagen y dieciocho clases agronómicas. Su desbalance es severo: la clase mayoritaria tiene
\num{6128} parcelas y la minoritaria \num{103}, una razón de unos \num{60}:\num{1}.

El \textbf{banco de réplica} es BreizhCrops \cite{russwurm2020breizhcrops}, series Sentinel-2 L2A
sobre Bretaña para la campaña de 2017, con nueve clases agronómicas y partición nativa por región
administrativa. Usamos las regiones \texttt{frh01} y \texttt{frh04} con un submuestreo
estratificado proporcional de \num{30000} parcelas cada una y semilla fija, \num{60000} en total.
Su desbalance es un orden de magnitud peor que el del banco primario: en el índice completo la
clase mayoritaria supera a la menor hasta por \num{19207}:\num{1}, y en nuestro submuestreo tres de
las nueve clases tienen dos, siete y \num{289} parcelas respectivamente. Los dos bancos se
describen con el mismo vector de \num{185} características temporales, derivado de diecisiete
índices espectrales, de modo que los dos experimentos difieren en dato y modelo pero no en
representación.

\begin{figure}[tbp]
\centering
\includegraphics[width=\textwidth]{soporte.pdf}
\caption{Soporte por clase en escala logarítmica. Las barras por debajo de \num{300} parcelas van
en el color de acento: son las clases que todo criterio de retirada quita primero. La cola del
banco de réplica no es solo más delgada, es degenerada —dos y siete parcelas—, y eso es lo que
hace posible el pasillo de la sección~\ref{sec:pasillo}.}
\label{fig:soporte}
\end{figure}

\subsection{Predictores}
\label{sec:predictores}

El protocolo opera sobre posteriores y es indiferente a cómo se produjeron. En el banco primario
tomamos los dos más fuertes de diez miembros disponibles, un transformador de visión para series de
imágenes satelitales \cite{tarasiou2023tsvit} y un modelo de gradiente potenciado sobre
incrustaciones aprendidas. En el banco de réplica entrenamos un clasificador de gradiente potenciado
\emph{dejando una región fuera}: la posterior de cada parcela viene de un modelo que nunca vio su
región, así que toda posterior que puntuamos es una predicción genuinamente fuera de región. Ese
clasificador alcanza F1-macro \num{0.4979} y exactitud \num{0.7421} en \texttt{frh01}, y
\num{0.4811} y \num{0.7257} en \texttt{frh04}.

\subsection{Bloques}
\label{sec:bloques}

Todo punto de operación se elige sobre datos disjuntos de los datos que lo miden. En el banco
primario los bloques son espaciales: una teselación H3 de resolución cinco, agrupada y separada por
un colchón de \SI{1}{\kilo\meter}, que da cinco bloques. En el banco de réplica el índice publicado
no trae coordenadas por parcela, así que la estructura espacial más fina disponible es la región
administrativa; usamos las dos regiones como los dos bloques. Es una limitación real y la tratamos
como tal: con dos bloques el remuestreo captura la variación \emph{dentro} de las regiones y no
\emph{entre} ellas, así que reportamos la diferencia por bloque de cada contraste y no solo un
intervalo agregado.

\subsection{Los dos mecanismos, y un tercer criterio}
\label{sec:mecanismos}

Sea $\mathcal{C}$ el conjunto completo de $C$ clases, $p(\cdot \mid x)$ la posterior del predictor
para la parcela $x$ y $\hat{y}(x) = \arg\max_{c \in \mathcal{C}} p(c \mid x)$ su predicción sin
restringir.

El \emph{recorte de leyenda} elige un subconjunto $\mathcal{L} \subseteq \mathcal{C}$ con
$|\mathcal{L}| = K$ sobre los bloques de entrenamiento, y en el bloque de evaluación entrega una
parcela cuando $\hat{y}(x) \in \mathcal{L}$, emitiendo $\arg\max_{c \in \mathcal{L}} p(c \mid x)$.
Probamos dos criterios para elegir $\mathcal{L}$: por F1 por clase sobre los bloques de
entrenamiento, y por soporte de clase sobre esos mismos bloques, el segundo porque es lo que se
reporta hacer en la práctica.

El \emph{rechazo por confianza} conserva $\mathcal{L} = \mathcal{C}$ y entrega las parcelas cuyo
$\max_c p(c \mid x)$ supera un umbral. El umbral es el cuantil de las confianzas del propio bloque
que iguala el número de parcelas que el mecanismo comparado entregó en ese mismo bloque, de modo
que los dos se comparan a igual cobertura.

Ningún mecanismo lee una etiqueta verdadera para decidir qué entrega. Vale la pena decirlo con
todas sus letras porque la implementación natural del recorte de leyenda sí la lee: entregar las
parcelas cuya etiqueta \emph{verdadera} cae en $\mathcal{L}$ es un oráculo, e infla la calidad
aparente del mecanismo. Nuestra regla de entrega usa el argmax sin restringir, que es observable en
inferencia.

\subsection{El estimando alineado}
\label{sec:estimando}

Dos mecanismos que prometen catálogos distintos no se pueden comparar por el promedio macro de cada
uno, porque los promedios se toman sobre denominadores distintos. Por eso puntuamos todo mecanismo
con

\begin{equation}
  \mathrm{F1}_{\text{alineado}} \;=\; \frac{1}{|\mathcal{L} \cap \mathcal{Y}|}
  \sum_{c \,\in\, \mathcal{L} \cap \mathcal{Y}} \mathrm{F1}_c ,
  \label{eq:alineado}
\end{equation}

\noindent donde $\mathcal{L}$ es la leyenda bajo prueba e $\mathcal{Y}$ las clases realmente
presentes entre las parcelas entregadas. Los dos mecanismos de un contraste se puntúan sobre la
\emph{misma} $\mathcal{L}$. Reportamos también el macro nativo de cada mecanismo sobre lo que
promete individualmente, pero solo como vista de producto; no es comparable entre mecanismos y
nunca lo usamos para un contraste.

\subsection{El control}
\label{sec:control}

El control que da título a este artículo cabe en una línea. Para cada leyenda $\mathcal{L}$ y cada
bloque, puntúese al predictor intacto sobre $\mathcal{L}$ con la ecuación~\ref{eq:alineado}
entregando \emph{todas} las parcelas. No se retira ninguna clase de la emisión ni se rechaza
ninguna parcela; solo cambia el denominador. Si ese control iguala o supera a un mecanismo, la
calidad la compró la elección del conjunto de clases y no el mecanismo.

\subsection{Inferencia}
\label{sec:inferencia}

Los contrastes usan un remuestreo pareado: se sortea un índice por bloque y se vuelven a puntuar
\emph{los dos} mecanismos sobre él, de modo que comparten el ruido de muestreo. Sorteamos
\num{1000} remuestreos y reportamos el intervalo percentil de la diferencia de medias del F1
alineado por bloque, junto con las diferencias por bloque. En el banco primario remuestreamos
además por parche de imagen en vez de por parcela, porque las parcelas de un mismo parche comparten
imagen y no son independientes. El criterio principal preregistrado es $K = \lceil C/2 \rceil$; los
demás valores de $K$ forman una familia exploratoria, corregida con el procedimiento descendente de
Holm y Bonferroni dentro de cada predictor y cada banco.
